\section{September 3, 2026}

Recall that if $\mu^*$ is an outer measure on $X$, a set $A\subseteq X$
is $\mu^*$-measurable if
\[
  \mu^*(E)=\mu^*(E\cap A)+\mu^*(E\cap A^c)
\]
for every $E\subseteq X$. By countable subadditivity, it is enough to
prove the inequality
\[
  \mu^*(E)\geq\mu^*(E\cap A)+\mu^*(E\cap A^c).
\]

\subsection{The Carath\'eodory construction}

\begin{theorem}[Carath\'eodory]\label{thm:caratheodory}
  Let $\mu^*$ be an outer measure on $X$, and set
  \[
    \mathcal M^*
    =\{A\subseteq X:A\text{ is $\mu^*$-measurable}\}.
  \]
  Then $\mathcal M^*$ is a $\sigma$-algebra, and
  \[
    \mu=\left.\mu^*\right|_{\mathcal M^*}
  \]
  is a complete measure on $\mathcal M^*$.
\end{theorem}

\begin{proof}
  We first show that $\mathcal M^*$ is an algebra. The empty set is
  measurable because
  \[
    \mu^*(E)=\mu^*(E\cap\varnothing)+\mu^*(E\cap X)
  \]
  for every $E\subseteq X$. The defining condition is symmetric in $A$
  and $A^c$, so
  \[
    A\in\mathcal M^*\quad\Longrightarrow\quad A^c\in\mathcal M^*.
  \]

  Now let $A,B\in\mathcal M^*$ and fix $E\subseteq X$. Applying the
  measurability criterion first to $A$ and then to $B$ gives
  \begin{align*}
    \mu^*(E)
    &=\mu^*(E\cap A)+\mu^*(E\cap A^c)\\
    &=\mu^*(E\cap A\cap B)+\mu^*(E\cap A\cap B^c)\\
    &\quad+\mu^*(E\cap A^c\cap B)
      +\mu^*(E\cap A^c\cap B^c).
  \end{align*}
  The first three sets on the last two lines have union
  $E\cap(A\cup B)$. Therefore countable subadditivity implies
  \[
    \mu^*(E)
    \geq\mu^*(E\cap(A\cup B))
      +\mu^*(E\cap(A\cup B)^c).
  \]
  Thus $A\cup B\in\mathcal M^*$, and $\mathcal M^*$ is an algebra.

  This argument also establishes finite additivity on disjoint
  measurable sets. Indeed, if $A,B\in\mathcal M^*$ are disjoint, apply
  the measurability of $A$ with $E=A\cup B$ to obtain
  \[
    \mu^*(A\cup B)
    =\mu^*((A\cup B)\cap A)
      +\mu^*((A\cup B)\cap A^c)
    =\mu^*(A)+\mu^*(B).
  \]

  To prove closure under countable unions, first let
  $A_1,A_2,\ldots\in\mathcal M^*$ be pairwise disjoint, and set
  \[
    B=\bigcup_{i=1}^{\infty}A_i,
    \qquad
    B_n=\bigcup_{i=1}^nA_i.
  \]
  Since $\mathcal M^*$ is an algebra, $B_n\in\mathcal M^*$. Repeatedly
  applying the measurability criterion and using the disjointness of the
  $A_i$ gives, for every $E\subseteq X$,
  \[
    \mu^*(E)
    =\sum_{i=1}^n\mu^*(E\cap A_i)+\mu^*(E\cap B_n^c).
  \]
  Because $B^c\subseteq B_n^c$, monotonicity yields
  \[
    \mu^*(E)
    \geq\sum_{i=1}^n\mu^*(E\cap A_i)+\mu^*(E\cap B^c).
  \]
  Letting $n\to\infty$ and then using countable subadditivity, we find
  \begin{align*}
    \mu^*(E)
    &\geq\sum_{i=1}^{\infty}\mu^*(E\cap A_i)
      +\mu^*(E\cap B^c)\\
    &\geq\mu^*\left(\bigcup_{i=1}^{\infty}(E\cap A_i)\right)
      +\mu^*(E\cap B^c)\\
    &=\mu^*(E\cap B)+\mu^*(E\cap B^c).
  \end{align*}
  Hence $B\in\mathcal M^*$. An arbitrary countable union can be
  disjointified by replacing $A_i$ with
  \[
    A_i\setminus\bigcup_{j=1}^{i-1}A_j.
  \]
  It follows that $\mathcal M^*$ is a $\sigma$-algebra.

  The same computation proves countable additivity. If the measurable
  sets $A_i$ are pairwise disjoint and $B=\bigcup_iA_i$, then countable
  subadditivity gives
  \[
    \mu^*(B)\leq\sum_{i=1}^{\infty}\mu^*(A_i).
  \]
  On the other hand, for every $n$,
  \[
    \mu^*(B)
    =\sum_{i=1}^n\mu^*(A_i)+\mu^*(B\setminus B_n)
    \geq\sum_{i=1}^n\mu^*(A_i).
  \]
  Letting $n\to\infty$ proves
  \[
    \mu^*(B)=\sum_{i=1}^{\infty}\mu^*(A_i).
  \]
  Thus $\mu=\left.\mu^*\right|_{\mathcal M^*}$ is a measure.

  Finally, suppose that $A\subseteq X$ and $\mu^*(A)=0$. For every
  $E\subseteq X$, subadditivity and monotonicity give
  \[
    \mu^*(E\cap A)+\mu^*(E\cap A^c)
    \leq0+\mu^*(E)=\mu^*(E).
  \]
  Hence $A\in\mathcal M^*$. In particular, if
  $N\in\mathcal M^*$ satisfies $\mu(N)=0$, then every subset $A\subseteq
  N$ has outer measure zero and therefore belongs to $\mathcal M^*$.
  This proves that $\mu$ is complete.
\end{proof}

\subsection{Premeasures}

The covering construction from the preceding lecture begins with a set
function $\rho$ defined on a family of elementary sets. The example
$\rho([a,b])=(b-a)^2$ showed that the resulting outer measure may fail to
agree with $\rho$, even on elementary sets. To prevent this collapse, we
require the initial set function to behave like a measure.

\begin{definition}[Premeasure]
  Let $\mathcal A$ be an algebra on $X$. A map
  \[
    \mu_0\colon\mathcal A\longrightarrow[0,\infty]
  \]
  is a \emph{premeasure} if
  \begin{enumerate}[label=(\roman*)]
    \item $\mu_0(\varnothing)=0$;
    \item if $A_1,A_2,\ldots\in\mathcal A$ are pairwise disjoint and
      $\bigcup_{i=1}^{\infty}A_i\in\mathcal A$, then
      \[
        \mu_0\left(\bigcup_{i=1}^{\infty}A_i\right)
        =\sum_{i=1}^{\infty}\mu_0(A_i).
      \]
  \end{enumerate}
\end{definition}

Thus a premeasure has the same additivity property as a measure, but its
domain is only required to be an algebra, rather than a $\sigma$-algebra.

\begin{example}[Half-open intervals]
  Let $X=\bb{R}$. Consider intervals of the forms
  \[
    (a,b],\qquad (-\infty,b],\qquad (a,\infty),
  \]
  together with $\varnothing$, where $a,b\in\bb{R}$ and $a<b$. The
  collection $\mathcal A$ of finite unions of such intervals is an
  algebra. Moreover,
  \[
    \mathcal M(\mathcal A)=\mathcal B_{\bb{R}}.
  \]
  This algebra will later be used to construct Lebesgue measure.
\end{example}

The next proposition identifies the hypotheses under which the outer
measure construction preserves the initial set function.

\begin{proposition}[Extension of a premeasure]
  Let $\mu_0$ be a premeasure on an algebra $\mathcal A$ over $X$. For
  every $E\subseteq X$, define
  \[
    \mu^*(E)
    =\inf\left\{
      \sum_{i=1}^{\infty}\mu_0(A_i):
      A_i\in\mathcal A,\quad
      E\subseteq\bigcup_{i=1}^{\infty}A_i
    \right\}.
  \]
  Then the following statements hold.
  \begin{enumerate}[label=(\roman*)]
    \item $\left.\mu^*\right|_{\mathcal A}=\mu_0$.
    \item Every member of $\mathcal A$ is $\mu^*$-measurable.
  \end{enumerate}
  Consequently, $\mathcal M(\mathcal A)\subseteq\mathcal M^*$, and the
  measure $\left.\mu^*\right|_{\mathcal M(\mathcal A)}$ extends
  $\mu_0$.
\end{proposition}

\begin{proof}
  Let $A\in\mathcal A$. The sequence
  $A,\varnothing,\varnothing,\ldots$ is a cover of $A$, so
  \[
    \mu^*(A)\leq\mu_0(A).
  \]
  Conversely, suppose that $A_i\in\mathcal A$ and
  $A\subseteq\bigcup_{i=1}^{\infty}A_i$. Define
  \[
    B_1=A\cap A_1,
    \qquad
    B_k=A\cap\left(A_k\setminus\bigcup_{i=1}^{k-1}A_i\right)
    \quad(k\geq2).
  \]
  Then the sets $B_k$ are pairwise disjoint members of $\mathcal A$,
  and
  \[
    A=\bigcup_{k=1}^{\infty}B_k.
  \]
  Since $B_k\subseteq A_k$, monotonicity of the premeasure gives
  \[
    \mu_0(A)
    =\sum_{k=1}^{\infty}\mu_0(B_k)
    \leq\sum_{k=1}^{\infty}\mu_0(A_k).
  \]
  Taking the infimum over all such covers yields
  $\mu_0(A)\leq\mu^*(A)$. Therefore $\mu^*(A)=\mu_0(A)$, proving part
  (i).

  For part (ii), fix $A\in\mathcal A$ and $E\subseteq X$. If
  $\mu^*(E)=\infty$, the desired measurability identity follows at once
  from countable subadditivity. Suppose instead that $\mu^*(E)<\infty$,
  and let $\varepsilon>0$. Choose $A_1,A_2,\ldots\in\mathcal A$ such
  that
  \[
    E\subseteq\bigcup_{i=1}^{\infty}A_i
    \qquad\text{and}\qquad
    \mu^*(E)+\varepsilon
    \geq\sum_{i=1}^{\infty}\mu_0(A_i).
  \]
  Since $\mu_0$ is finitely additive,
  \begin{align*}
    \sum_{i=1}^{\infty}\mu_0(A_i)
    &=\sum_{i=1}^{\infty}
      \bigl(\mu_0(A_i\cap A)+\mu_0(A_i\cap A^c)\bigr)\\
    &\geq\mu^*(E\cap A)+\mu^*(E\cap A^c).
  \end{align*}
  Hence
  \[
    \mu^*(E)+\varepsilon
    \geq\mu^*(E\cap A)+\mu^*(E\cap A^c).
  \]
  Letting $\varepsilon\downarrow0$ proves that $A$ is
  $\mu^*$-measurable.

  By Theorem~\ref{thm:caratheodory}, $\mathcal M^*$ is a
  $\sigma$-algebra. Since it contains $\mathcal A$, it also contains
  $\mathcal M(\mathcal A)$. Part (i) then shows that the restriction of
  $\mu^*$ to $\mathcal M(\mathcal A)$ extends $\mu_0$.
\end{proof}
